% Bản nháp chương 7 theo dàn bài mới, viết với bộ mẫu ở style/mau/.
% ===========================================================================
\chapter{Lý thuyết so với thực nghiệm}
\label{sec:theory}

Ba cơ chế ở chương~\ref{sec:mechanisms} đều rút ra từ số đo. Lý thuyết tối ưu hóa
lồi cũng nói về đúng những đại lượng ấy, và nói trước khi chạy bất kỳ thí nghiệm
nào. Chương này đặt hai bên cạnh nhau để xem cận lý thuyết mô tả đúng chỗ nào,
lỏng ở chỗ nào, và dùng được vào việc gì.

Bài toán ở đây cho phép làm phép đối chiếu ấy chặt hơn thường lệ. Vì $\fstar$
tính được trước từ \eqref{eq:closed-form} nên hệ số co rút quan sát được đo bằng
một cái thước tuyệt đối, chứ không phải ước lượng từ giá trị tốt nhất mà thuật
toán tình cờ đạt tới.

\section{Cận lý thuyết}
\label{sec:theory-bound}

Với hàm bậc hai lồi mạnh và độ dài bước $t = 2/(L+\mu)$,
%
\begin{equation}
  \label{eq:gd-rate}
  f(w_k) - \fstar \le \Bigl( \frac{\kappa - 1}{\kappa + 1} \Bigr)^{2k}
                      \bigl( f(w_0) - \fstar \bigr) ,
\end{equation}
%
còn phương pháp tăng tốc có hệ số co rút cỡ $1 - \kappa^{-1/2}$~\cite{nesterov2018lectures}.

Đọc \eqref{eq:gd-rate} thành lời thì mỗi vòng lặp nhân sai số với một hệ số cố
định nhỏ hơn 1, và hệ số ấy tiến về 1 khi $\kappa$ lớn. Nó đến thẳng từ cơ chế
thứ nhất ở mục~\ref{sec:mech-step}: sai số theo hướng riêng ứng với trị riêng
$\lambda_i$ nhân với $\abs{1 - t\lambda_i}$ mỗi vòng, và với $t = 2/(L+\mu)$ thì
giá trị lớn nhất của $\abs{1 - t\lambda_i}$ trên khoảng $[\mu, L]$ đạt được ở cả
hai đầu và bằng $(\kappa-1)/(\kappa+1)$. Cận cho $f$ vì thế là bình phương của
con số đó, tức \num{0.985202} với $\kappa = \num{268.31}$.

Điểm cần giữ trong đầu khi đọc phần sau: \eqref{eq:gd-rate} là cận cho trường hợp
xấu nhất trong lớp hàm lồi mạnh có $\kappa$ cho trước, không phải dự đoán cho một
hàm cụ thể. Một bài toán riêng lẻ được phép nhanh hơn cận, và câu hỏi đáng hỏi là
nhanh hơn bao nhiêu và vì sao.

\section{Hệ số co rút ở phần đuôi}
\label{sec:theory-tail}

\resultfig{theory_vs_practice}
  {Đường hội tụ thực nghiệm đặt cạnh hai cận lý thuyết. Hình này chỉ có bản theo
   số vòng lặp, vì cận lý thuyết phát biểu theo vòng lặp và không có tương ứng
   trên trục thời gian.}
  {fig:theory}

Nhóm ước lượng hệ số co rút thực nghiệm bằng cách khớp đường thẳng vào phần tuyến
tính ở đuôi của đồ thị $\log_{10}\bigl(f(w_k) - \fstar\bigr)$ theo $k$, rồi đối
chiếu với cận trong bảng~\ref{tab:rates}.

\begin{table}[tp]
  \centering
  \caption{Hệ số co rút mỗi vòng lặp trên $f$, quan sát được ở phần đuôi so với
    cận lý thuyết.}
  \label{tab:rates}
  \begin{tabular}{lrr}
    \toprule
    Phương pháp & Quan sát & Cận lý thuyết \\
    \midrule
    GD, $t = 2/(L+\mu)$                     & \num{0.98520} & \num{0.98520} \\
    GD, $t = 2/L$                           & \num{0.99999} & --            \\
    AGD, $\beta$ hằng tính từ $\kappa$      & \num{0.78894} & \num{0.93895} \\
    AGD, $\beta_k$ tăng dần + khởi động lại & \num{0.65455} & \num{0.93895} \\
    \bottomrule
  \end{tabular}
\end{table}

Với gradient descent ở độ dài bước mà \eqref{eq:gd-rate} nhắm tới, hệ số quan sát
được là \num{0.98520} và trùng với cận tới năm chữ số thập phân
(bảng~\ref{tab:rates}, hình~\ref{fig:theory}). Cận không chỉ đúng mà còn chặt, vì
ở giai đoạn đuôi đúng thành phần ứng với trị riêng nhỏ nhất chi phối sai số, và
tốc độ tắt của thành phần đó chính là $(\kappa-1)/(\kappa+1)$.

Dòng thứ hai của bảng là phép thử ngược cho cùng một cơ chế. Với $t = 2/L$, hệ số
quan sát được là \num{0.99999}, tức thực tế bằng 1, đúng như dự đoán từ
mục~\ref{sec:mech-step} rằng thành phần ứng với $\lambda_{\max} = L$ đứng yên ở
độ dài bước ấy. Cận \eqref{eq:gd-rate} không áp cho cấu hình này vì nó phát biểu
cho $t = 2/(L+\mu)$, nên ô tương ứng để trống.

Với phương pháp tăng tốc, hệ số quan sát được là \num{0.65455} so với cận
\num{0.93895}, tức tốt hơn cận rất nhiều nhưng không mâu thuẫn với nó, vì cận là
cận trên. Điều đáng chú ý hơn nằm ở hai dòng cuối bảng: công thức momentum hằng
tính từ $\kappa$, vốn là công thức tối ưu về mặt lý thuyết, cho hệ số
\num{0.78894}, kém hơn công thức tăng dần kèm khởi động lại. Cấu hình được lý
thuyết chỉ định không phải cấu hình chạy nhanh nhất trên bài toán này, và mục sau
giải thích vì sao.

\section{Vì sao cận nói rất ít về giai đoạn đầu}
\label{sec:theory-early}

Chênh lệch lớn nhất không nằm ở hệ số co rút mà nằm ở số vòng lặp. Gradient
descent đạt $10^{-6}$ sau 20 vòng lặp (bảng~\ref{tab:gd-fixed}), trong khi ước
lượng từ \eqref{eq:gd-rate} cho khoảng $\tfrac{1}{2}\kappa \ln(10^{6}) \approx
1850$ vòng, tức lệch 93 lần.

Nguyên nhân đo được, và nó nằm ở phổ của Hessian. Cận \eqref{eq:gd-rate} là chặt
khi khối lượng sai số ban đầu trải đều trên toàn khoảng $[\mu, L]$, vì khi ấy
thành phần chậm nhất chiếm phần đáng kể ngay từ đầu. Trên bài toán này thì không:
trong 280 trị riêng của Hessian, 269 trị nằm trên $L/10$ và chỉ 8 trị nằm dưới
$10\mu$, còn trị riêng trung vị bằng \num{0.377} lần $L$. Nói cách khác, phổ dồn
về phía đầu lớn, và tám hướng phẳng mà $\lambda$ tạo ra ở
\eqref{eq:mu-equals-lambda} chỉ chiếm một phần rất nhỏ của sai số ban đầu.

Hệ quả là hai giai đoạn tách hẳn nhau. Trong vài chục vòng đầu, sai số do 269
hướng cong chi phối và chúng tắt với hệ số cỡ $\abs{1 - tL}$ chứ không phải
$(\kappa-1)/(\kappa+1)$, nên tốc độ nhanh hơn cận nhiều bậc. Chỉ khi các hướng ấy
đã tắt gần hết thì tám hướng phẳng mới lộ ra, và từ đó trở đi cận mô tả đúng tới
năm chữ số thập phân như bảng~\ref{tab:rates} cho thấy.

Cùng lý do đó giải thích hai dòng AGD. Công thức momentum hằng tính từ $\kappa$
được thiết kế để cân bằng đúng trường hợp xấu nhất, tức trường hợp phổ trải đều,
nên trên một phổ dồn về đầu lớn nó hãm phần nhanh lại nhiều hơn mức cần. Công
thức tăng dần không giả định gì về $\kappa$ nên nó không mắc lỗi ấy, và khởi động
lại cắt nốt phần dao động mà momentum tăng dần sinh ra ở đuôi. Trên một bài toán
có phổ trải đều giữa $\mu$ và $L$, cả hai chênh lệch này đều thu hẹp lại.

\section{Cận dùng được vào việc gì}
\label{sec:theory-use}

Đối chiếu ở trên cho ba kết luận thực dụng.

Cận \eqref{eq:gd-rate} dự đoán đúng hành vi ở phần đuôi và có thể dùng để trả lời
câu hỏi chạy thêm bao lâu thì cắt thêm được một bậc sai số, một khi thuật toán đã
vào vùng tuyến tính. Ngược lại, nó không dùng được để lập ngân sách vòng lặp từ
đầu: con số 1850 vòng lặp lớn hơn thực tế 93 lần, và lập kế hoạch theo nó thì
đúng nhưng lãng phí gần hai bậc.

Cận cũng không dùng được để chọn cấu hình. Momentum hằng tính từ $\kappa$ là lựa
chọn tối ưu theo lý thuyết nhưng chậm hơn momentum tăng dần trên bài toán này,
nên việc chọn tham số vẫn phải dựa vào lưới thí nghiệm ở
chương~\ref{sec:mechanisms} chứ không dựa vào công thức.

Cuối cùng, chỗ cận thực sự có ích là làm phép kiểm tra tính đúng đắn. Một hệ số
co rút quan sát được nằm trên cận, chẳng hạn lớn hơn \num{0.98520} với cấu hình ở
dòng đầu bảng~\ref{tab:rates}, là dấu hiệu của lỗi cài đặt chứ không phải của một
bài toán khó, vì cận đúng cho mọi hàm trong lớp. Nhóm dùng đúng phép kiểm ấy khi
truy lỗi momentum ở mục~\ref{sec:mech-momentum}.

Mọi con số trong chương này đo ở đúng một giá trị $\lambda$, tức ở đúng một
$\kappa$. Vì $\kappa$ có mặt trong cả hai cận, chương~\ref{sec:lambda} kiểm xem
mức độ khớp giữa lý thuyết và thực nghiệm có giữ nguyên khi đổi $\kappa$ hay
không.
